\begin{frame}
    \frametitle{\insertsectionhead}
    \begin{block}{Kruskal算法}
        \begin{enumerate}
            \item 将边按权重从小到大排序；初始时每个顶点各自构成一个连通分量
            \item 依次检查边：若两端点属于不同分量，则选中该边并合并两个分量；
                  否则跳过，避免形成环
            \item 选中$n-1$条边时结束；可使用并查集维护连通分量
        \end{enumerate}
    \end{block}
    \begin{alertblock}{适用范围}
        对连通无向带权图求最小生成树；图不连通时得到最小生成森林。
    \end{alertblock}
\end{frame}

\begin{fragile}
    \frametitle{\insertsectionhead}
    \begin{figure}
        \centering
        \begin{tikzpicture}[ >=stealth, thick, black!50, %
                list item/.style={draw=gray, circle, thick}, %
            ]
            \node (a) [terminal] at (0,0) {$A$};
            \node (b) [terminal] at ($(a.center)+(18:2)$) {$B$};
            \node (c) [terminal] at ($(a.center)+(90:2)$) {$C$};
            \node (d) [terminal] at ($(a.center)+(162:2)$) {$D$};
            \node (e) [terminal] at ($(a.center)+(234:2)$) {$E$};
            \node (f) [terminal] at ($(a.center)+(306:2)$) {$F$};
            \only<-6>{\draw (a) to node [pos=0.5,] {$26$} (b);}
            \only<7->{\draw[red,very thick] (a) to node [pos=0.5,] {$26$} (b);}
            \only<-3>{\draw (a) to node [pos=0.5,] {$19$} (d);}
            \only<4->{\draw[red,very thick] (a) to node [pos=0.5,] {$19$} (d);}
            \only<-4>{\draw (a) to node [pos=0.5,] {$25$} (e);}
            \only<5->{\draw[red,very thick] (a) to node [pos=0.5,] {$25$} (e);}
            \only<-5,7->{\draw (a) to node [pos=0.5,] {$25$} (f);}
            \only<6>{\draw[red,dashed] (a) to node [pos=0.5,] {$25$} (f);}
            \only<1>{\draw (b) to node [pos=0.5,] {$12$} (c);}
            \only<2->{\draw[red,very thick] (b) to node [pos=0.5,] {$12$} (c);}
            \draw (c) to node [pos=0.5,] {$34$} (d);
            \draw (d) to node [pos=0.5,] {$46$} (e);
            \only<-2>{\draw (e) to node [pos=0.5,] {$17$} (f);}
            \only<3->{\draw[red,very thick] (e) to node [pos=0.5,] {$17$} (f);}
            \draw (f) to node [pos=0.5,] {$38$} (b);
        \end{tikzpicture}
        \caption{Kruskal算法示例过程}
        \label{fig:demo_mst_kruskal}
    \end{figure}
    \begin{center}
        \small
        粗实线：已选边；红色虚线：当前检查但跳过的边。\\
        \only<1>{初始时各顶点分别为一个分量；按边权递增检查。}
        \only<2>{选中$BC(12)$，合并$B$与$C$所在分量。}
        \only<3>{选中$EF(17)$，合并$E$与$F$所在分量。}
        \only<4>{选中$AD(19)$，合并$A$与$D$所在分量。}
        \only<5>{$AE$与$AF$同权，先选$AE(25)$，合并两个分量。}
        \only<6>{跳过$AF(25)$：$A$与$F$已连通，加入会形成环。}
        \only<7>{选中$AB(26)$；已有$5$条边，总权重为$99$。}
    \end{center}
\end{fragile}

\begin{frame}
    \frametitle{\insertsectionhead}
    \begin{exampleblock}{前述网线铺设示例}
        \begin{center}
            \begin{tabular}{c|c|l}
                \toprule
                边 & 权重 & 处理 \\
                \midrule
                $BC$ & $12$ & 选中 \\
                $EF$ & $17$ & 选中 \\
                $AD$ & $19$ & 选中 \\
                $AE$ & $25$ & 选中 \\
                $AF$ & $25$ & 跳过：与$AE,EF$形成环 \\
                $AB$ & $26$ & 选中，已有$5$条边 \\
                \bottomrule
            \end{tabular}
        \end{center}
        权重之和为$12+17+19+25+26=99$；同权边顺序不同时结果可能不同。
    \end{exampleblock}
\end{frame}

\begin{frame}
    \frametitle{\insertsectionhead}
    \begin{table}
        \centering
        \small
        \caption{Prim与Kruskal的比较（$n$个顶点，$e$条边）}
        \begin{tabular}{p{0.20\textwidth}p{0.35\textwidth}p{0.35\textwidth}}
            \toprule
            & Prim & Kruskal \\
            \midrule
            核心思路 & 从一个顶点扩展生成树 & 按边权合并不同分量 \\
            常见实现 & 邻接矩阵；或邻接表加索引堆 & 边排序加并查集 \\
            时间 & 矩阵$O(n^2)$；堆$O((n+e)\log n)$ & $O(n+e\log e)$ \\
            存储 & 矩阵$O(n^2)$；邻接表$O(n+e)$ & 边集及并查集$O(n+e)$ \\
            \bottomrule
        \end{tabular}
    \end{table}
    \small 存储包括图或边集本身；矩阵适合稠密图，边排序常用于稀疏图。
\end{frame}
